% ==============================================================================
% Main paper: Annotation Pipeline
%==============================================================================

\section{Building Artifact-aware Agent-as-a-judge}
\label{sec:judge}

To make \bench to be fully utilized and uncover hidden failure patterns behind the trajectories, a systematic and comprehensive analysis method is required. However, many failure patterns leave no
trace in the report at all: the agent may claim a result its own code does
not produce or describe a method its logs show it never ran. Detecting these failures
requires comparing the manuscript against the full set of artifacts the agent
produced. We
therefore build an annotation pipeline in three steps: human annotation to
develop and calibrate the failure taxonomy (\S\ref{sec:taxonomy}), an
automated Agent-as-a-Judge to scale annotation to the full 800-trajectory
corpus, and a validation study measuring the judge's agreement with human
labels.

\subsection{Human Annotation}
\label{sec:judge:human}

Our failure taxonomy was developed inductively from expert examination of
agent trajectories, following a grounded-theory process\citep{glaser1967discovery}. Rather than
starting from a predefined checklist, experts examined complete
trajectories independently---execution logs, delivered code, final reports, and data files---and recorded whatever failure behaviors they observed. Observed
patterns were iteratively grouped, split, and refined through constant
comparative analysis until further annotation yielded no new failure modes
(theoretical saturation). We provide a detailed description and analysis of
the resulting taxonomy in \S~\ref{sec:taxonomy}.

To validate that the taxonomy can be applied consistently, we conduct
inter-annotator agreement (IAA) studies on a stratified sample of
$50$ trajectories drawn from all $800$ trajectories.
Three experts independently label each sampled trajectory against
the full taxonomy. Taxonomy then is iteratively sharpened and refined---adding, merging, or
clarifying categories---until three experts reach consensus. We conduct $5$
rounds of IAA, achieving $\kappa = 0.85$ (Cohen's Kappa) in the final
round. Because annotation difficulty varies across failure types---failures
that leave concrete artifacts (e.g.\ a code file contradicting the report)
are easier to adjudicate than failures requiring metacognitive
judgment---agreement is not uniform across the taxonomy, and we treat
patterns in the Cognitive Depth \& Adaptability pillar as lower-confidence
throughout (Appendix~\ref{app:evaluator:labeling}).

\begin{table}[t]
\centering
\small
\caption{Agreement with human expert annotation on the 50 validation
trajectories. The LLM-as-a-Judge baseline receives only the transcript in
a single call; the Agent-as-a-Judge receives the full evidence package and
operates under the structured rubric. \emph{Pattern} measures per-pattern
hit/miss agreement across the 45 \taxonomy patterns; \emph{Taxonomy
Categorization} measures agreement at the root-cause pillar level.}
\label{tab:judge-performance}
\begin{tabular}{llccccc}
\toprule
Method & Level & Accuracy & Precision & Recall & F1 & Cohen's $\kappa$ \\
\midrule
\multirow{2}{*}{LLM-as-a-Judge (claude-opus-5)}
 & Pattern       & 84.6 & 70.2 & 63.5 & 66.7 & 0.53 \\
 & Taxonomy Categorization & 89.3 & 78.8 & 72.1 & 75.3 & 0.62 \\
\midrule
\multirow{2}{*}{Agent-as-a-Judge}
 & Pattern       & 92.1 & 85.4 & 80.7 & 83.0 & 0.75 \\
 & Taxonomy Categorization & 95.3 & 91.0 & 87.2 & 89.1 & 0.83 \\
\bottomrule
\end{tabular}
\end{table}

\subsection{Agent-as-a-Judge}
\label{sec:judge:agent}

Manually annotating 800 trajectories at this depth is prohibitively
expensive and time-consuming. To scale annotation we develop an
\textbf{artifact-aware} Agent-as-a-Judge: an autonomous agent that
analyzes a trajectory for failure patterns across all lifecycle stages and
categorizes them according to the \taxonomy. Unlike a single LLM-as-a-judge
call on the transcript, the agent judge can execute code and navigate the
rollout's full artifact set so
failures invisible in the report alone become detectable.
 
The judge receives the rollout's complete evidence package:  code, execution logs, final reports, and data files (Appendix~\ref{app:evaluator:evidence}). Its prompt
enforces a fixed rubric aligned with the lifecycle stages of the taxonomy,
requiring the judge to cover each stage in a dedicated section and support
every identified issue with verifiable evidence---a log line number, file
name, code identifier, or exact numeric value. A set of nine \emph{iron
rules}, distilled from earlier annotation rounds, guard against common
mis-judgments (Appendix~\ref{app:evaluator:rubric}). An automated quality
checker enforces coverage, depth, and anchor density before accepting a
document; documents that fail are regenerated with gate-specific feedback
in a self-healing loop (Appendix~\ref{app:evaluator:loop}). Accepted
analyses are then mapped to taxonomy pattern IDs in a labeling pass,
producing the failure counts aggregated in
Figure~\ref{fig:failure-attribution}. Full details of the judge harness,
rubric, checker, and labeling protocol are in
Appendix~\ref{app:agent-as-a-judge}.

\subsection{Validation Against Human Annotation}
\label{sec:judge:validation}

We validate the Agent-as-a-Judge against human expert annotations on the 50 \add{human-labeled }%
trajectories of \S\ref{sec:judge:human}.
Each trajectory is annotated independently by the
judge, and we measure agreement at the failure pattern and taxonomy categorization level against human experts. For level pattern measurement, human experts and LLM are involved to compare and judge the results.
Table~\ref{tab:judge-performance} reports the results.

The artifact-aware Agent-as-a-Judge substantially outperforms the single-call LLM-as-a-judge at both granularities, with the largest gains in recall ($+17.2$ at pattern level), indicating that artifact access is required in practice for detecting failures invisible in the transcript.

